% 20260526_Assingment_Algebra_IIA
\documentclass[dvipdfmx]{jsreport}
\usepackage{soupreamble}
\title{代数学特論IIA 第1回レポート課題}
\author{M6012 川村聡一郎}
\date{}

\begin{document}
\maketitle

\begin{tcolorbox}
	$\mu$を正実数とする．このとき次を示せ．
\begin{enumerate}
	\item $\sum_{k= 1}^\infty \frac{1}{k^\mu}$は$\mu> 1$のとき収束する．
	\item $\sum_{k= 1}^\infty \frac{1}{k^\mu}$は$\mu\le 1$のとき発散する．
\end{enumerate}
\end{tcolorbox}

\begin{souanswer}
\begin{enumerate}
	\item 関数$f(x)= \frac{1}{x^\mu}$を考える．$\mu> 0$であるから，区間$[1, \infty)$で$f(z)$は正かつ連続で単調減少．
\begin{equation*}
	\int_1^M f(x)dx= \int_1^M \frac{1}{x^\mu} dx= \left[\frac{x^{1- \mu}}{1- \mu}\right]_1^M= \frac{M^{1- \mu}}{1- \mu}
\end{equation*}
	ここで$1- \mu< 0\ (\because \mu> 1)$．したがって$M\to \infty$で$M^{1- \mu}\to 0$．よって$\sum_{k= 1}^\infty \frac{1}{k^\mu}$は$\mu> 1$で収束．
	
	\item 
\begin{align*}
	\sum_{k= 1}^\infty \frac{1}{k^\mu}
	&= \frac{1}{1^\mu}+ \frac{1}{2^\mu}+ \frac{1}{3^\mu}+ \frac{1}{4^\mu}+ \cdots\\
	&> \frac{1}{1^\mu}+ \frac{1}{2^\mu}+ \frac{1}{2^\mu}+ \frac{1}{4^\mu}+ \cdots\\
	&= \frac{1}{2^{0\cdot \mu}}+ \frac{1}{2^{1\cdot \mu}}\times 2+ \frac{1}{2^{2\cdot \mu}}\times 4+ \cdots\\
\end{align*}
	これは初項1，公比$2^{1- \mu}(> 0)$の等比級数であり，この級数の和は$\frac{1\times\left(1- \left(\frac{2}{2^\mu}\right)^n\right)}{1- \frac{2}{2^\mu}}\to \infty\ (n\to \infty)$となるから，$\sum_{k= 1}^\infty \frac{1}{k^\mu}$は$\mu\le 1$で発散する．
\end{enumerate}	
\end{souanswer}


\begin{tcolorbox}
	$\omega_1, \omega_2\in \C$に対し，次が成立することを示せ．
\begin{enumerate}
	\item $\omega_1\ne 0\Longleftrightarrow \omega_1$は$\C$上一次独立
	\item $\omega_1$と$\omega_2$が$\R$上一次独立$\Longleftrightarrow \frac{\omega_1}{\omega_2}\not\in \R$
\end{enumerate}
\end{tcolorbox}

\begin{souanswer}
\begin{enumerate}
	\item 
\begin{enumerate}
	\item[$\Rightarrow)$] 一次独立の定義より$c_1\in \C$に対し，$c_1\omega_1= 0\Rightarrow c_1= 0$
	$\omega_1\ne 0$と仮定すると$c_1\omega_1= 0$ならば$c_1= 0$．一次独立の定義を満たし，$\omega_1$は$\C$上で一次独立であるといえる．
	\item [$\Leftarrow)$] 対偶を示す．すなわち$\omega_1= 0$と仮定し，一次独立の定義を満たさないことを示す．\\
	$\omega_1= 0$と仮定する．このとき$c_1\omega_1= 0$をみたす係数$c_1\in \C$が存在する．$\omega_1= 0$であったから$c_1$は$c_1= 0$に限定されず，任意の複素数で成り立つ．これは一次独立の定義を満たさないため矛盾．したがって$\omega_1\ne 0$でなければならない
\end{enumerate}
\begin{enumerate}

	\item[$\Rightarrow)$] 対偶を示す．
	一次従属の定義より$c_1, c_2\in \R$(ただし少なくとも一方は0ではない)が存在して，$c_1\omega_1+ c_2\omega_2= 0$とかける．
\begin{enumerate_roman}
	\item $c_1= 0$のとき：$c_1\omega_1+ c_2\omega_2= 0+ c_2\omega_2= 0$．$\omega_2\ne 0$であったからこの等式を満たすためには$c_2= 0$とならなければならないが，これは一次従属の定義をみたさない．

    \item $c_1\ne 0$のとき：$\frac{\omega_1}{\omega_2}= -\frac{c_2}{c_1}$と変形できる．$c_1, c_2\in \R$であるから$-\frac{c_2}{c_1}\in \R$となり，$\frac{\omega_1}{\omega_2}\in \R$が示された．
\end{enumerate_roman}
	\item[$\Leftarrow)$] 対偶を示す．ある実数$k\in \R$を用いて$\frac{\omega_1}{\omega_2}= k$とかけると仮定する．両辺に$\omega_2$をかけて整理すると$\omega_1- k\omega_2= 0$．これは$\omega_1, \omega_2$の係数が実数かつ少なくとも一方は0ではないため，一次従属の定義をみたす
\end{enumerate}
\end{enumerate} 
\end{souanswer}









\end{document}
